%% main.tex -- arXiv preprint
%% Verifiable Append-Only Governance: Kernel-Proven Acyclicity, FIFO,
%%   and Monotonicity for AI Decision Receipts
%% Author: Stephen P. Lutar Jr. (ORCID 0009-0001-0110-4173), SZL Holdings
%% License: CC-BY-4.0
%% arXiv target: cs.LO (primary); cross-list cs.CR, cs.DC, cs.LG
%% Doctrine v11: honest tiering; locked-proven = exactly 8.

\documentclass[11pt,letterpaper]{article}
\usepackage[margin=1in]{geometry}
\usepackage{amsmath,amssymb,amsthm}
\usepackage{booktabs}
\usepackage{array}
\usepackage{listings}
\usepackage{xcolor}
\usepackage{placeins}
\usepackage{multicol}
\usepackage[colorlinks=true,linkcolor=blue,citecolor=blue,urlcolor=blue]{hyperref}
\setlength{\emergencystretch}{3em}
\Urlmuskip=0mu plus 1mu

\lstset{basicstyle=\ttfamily\small,breaklines=true,frame=single,
  columns=fullflexible,keepspaces=true,xleftmargin=1em,framexleftmargin=1em}

\newtheorem{theorem}{Theorem}[section]
\newtheorem{corollary}[theorem]{Corollary}
\newtheorem{lemma}[theorem]{Lemma}
\theoremstyle{definition}
\newtheorem{definition}[theorem]{Definition}
\newtheorem{remark}[theorem]{Remark}
\newtheorem{conjecture}[theorem]{Conjecture}

\newcommand{\Lam}{\Lambda}
\newcommand{\kernelaxioms}{\{\texttt{propext},\,\texttt{funext},\,\texttt{Classical.choice},\,\texttt{Quot.sound}\}}

\title{\textbf{Verifiable Append-Only Governance:}\\
  Kernel-Proven Acyclicity, FIFO Reception, and Emit Monotonicity\\
  for AI Decision Receipts}

\author{Stephen P. Lutar Jr.\thanks{SZL Holdings, Inc. ORCID
  \href{https://orcid.org/0009-0001-0110-4173}{0009-0001-0110-4173}.
  Correspondence: \texttt{stephenlutar2@gmail.com}. License CC-BY-4.0.
  Related program record (not this paper's DOI):
  \href{https://doi.org/10.5281/zenodo.19944926}{10.5281/zenodo.19944926}.
  Lean development: \href{https://github.com/szl-holdings/lutar-lean}{github.com/szl-holdings/lutar-lean}.}}

\date{July 2026}

\begin{document}
\maketitle

\begin{abstract}
Agentic AI is \emph{post-deterministic}: the same prompt and tools may yield
different action sequences, so accountability cannot rest on re-deriving a
deterministic trace. It must instead rest on a \emph{record} whose structural
invariants are themselves verifiable. We give a machine-checked foundation for
three such invariants over an AI decision-receipt substrate, with proof terms
checked by Lean~4 at the named snapshot: (F4) the modeled receipt directed
acyclic graph (DAG) is acyclic and remains acyclic under append; (F7) draining
a modeled \emph{Chaski} batch preserves first-in-first-out (FIFO) order, with a
positional witness that the $i$-th received message equals the $i$-th sent
message; and (F22) the constructed range-based sequence log extends by one and
is strictly monotone by position. These are model-level results; durable broker
ordering and storage overwrite prevention remain deployment refinement
obligations.
These three results, kernel-verified and non-vacuous as of 2026-06-10, raise
SZL Holdings' \emph{locked-proven} formula count from five to exactly eight,
$\{F1,F4,F7,F11,F12,F18,F19,F22\}$ --- a count itself certified by the kernel
($\texttt{locked\_count\_eight} : \texttt{lockedDisclosed.length} = 8$, by
\texttt{decide}). We position the contribution honestly: the underlying graph
and queue facts are classical, but their assembly into a single,
audit-oriented, axiom-disclosed \emph{governance receipt} layer for
post-deterministic agents, with a published disclosure-name registry, a
reproducible per-declaration \texttt{\#print axioms} audit protocol, and a hard
wall between the locked eight and a $\sim$185-theorem experimental tier, is,
to our knowledge, new. We are explicit
about what is \emph{not} proven: cryptographic binding is axiom-gated under a
declared collision-resistance idealization, Conjecture~1's unconditional
$\Lambda$-uniqueness statement is disproved as stated while a weaker-condition
question remains open, and Khipu Byzantine safety remains Conjecture~2.
\end{abstract}

\noindent\textbf{Keywords:} append-only ledger, decision receipts, directed
acyclic graph, FIFO ordering, formal verification, Lean~4, Mathlib, AI
governance, supply-chain provenance, honest claims.

\section{Introduction}\label{sec:intro}

Classical software governance assumes determinism: an auditor who wants to know
what a program did can in principle re-run it on the recorded inputs and
recover the recorded outputs. Agentic AI dissolves this assumption. A
large-model agent that plans, calls tools, and revises its plan is
\emph{post-deterministic}: its control flow is sampled, its environment is
non-stationary, and the ``same'' task may produce different yet individually
defensible action sequences~\cite{brundage2020}. Re-derivation no longer
establishes accountability.

The auditing community has converged on a sharper requirement. It is no longer
sufficient to answer ``who did what''; one must answer \emph{under what
authority} an action occurred, with a record an auditor cannot dispute, and the
strongest framings demand evidence captured \emph{outside the agent's own
process boundary}~\cite{isaca2025,armo2025}. This shifts the burden from the
agent's reasoning (irreducibly non-deterministic) to the agent's
\emph{record} --- and a record is only as trustworthy as the structural
invariants one can prove about it.

This paper isolates and machine-verifies three such invariants over an AI
decision-receipt substrate. Each is a property that an auditor would want to
\emph{assume} but should instead be able to \emph{check}:

\begin{enumerate}
\item \textbf{Acyclicity of the receipt DAG (F4).} Receipts reference prior
  receipts, forming a DAG. If a cycle could form, ``provenance'' would be
  circular and an attacker could fabricate a self-justifying history. We prove
  the receipt DAG is acyclic and \emph{stays} acyclic under append.
\item \textbf{FIFO reception ordering (F7).} Inter-organ messages (the
  \emph{Chaski} channel) should be delivered in send order for causal audit
  reconstruction. We prove that the modeled batch drain returns send order,
  positionally; a durable broker must separately refine this model.
\item \textbf{Emit monotonicity (F22).} The constructed range log extends its
  prior sequence and its values strictly increase with position. Preventing
  overwrite in a deployed store remains a refinement obligation.
\end{enumerate}

\paragraph{Contribution and honest scope.} We do not claim the underlying
mathematics is new. Topological acyclicity of a backward-edge DAG, FIFO
queue order-preservation, and strict monotonicity of an index sequence are
elementary. Our contribution is fourfold and we state it precisely:
(i) the \emph{assembly} of exactly these three invariants into a single,
audit-oriented governance receipt layer for post-deterministic agents;
(ii) their \emph{machine-checked, non-vacuous} formalization in Lean~4 over
explicit data models (an edge-set DAG, a FIFO queue, and a range-based emit
log), together with a published disclosure-name registry and a protocol for
capturing per-declaration \texttt{\#print axioms} output;
(iii) a kernel-certified \emph{count} discipline --- the locked-proven set is
exactly eight, and ``eight'' is itself a theorem (\S\ref{sec:count}); and
(iv) the honesty architecture that walls the locked eight off from a
$\sim$185-theorem experimental tier and that flags two prior versions of F4/F7
as \emph{vacuous} and replaces them with genuine statements
(\S\ref{sec:nonvacuity}). The last point is, we believe, the most transferable:
we show by example how to detect and repair a vacuous formalization rather than
ship it.

\paragraph{What this paper does not claim.} Cryptographic binding of the
hash-chain (you cannot alter a past entry without breaking the Merkle root) is
\emph{axiom-gated} under a declared, disclosed collision-resistance
idealization; it is not a proof of cryptographic hardness (\S\ref{sec:binding}).
The trust aggregator $\Lambda$ (the equal-weight geometric mean used elsewhere
in the SZL substrate) had an unconditional uniqueness statement labeled
\textbf{Conjecture~1}; that statement is machine-checked \emph{false}, a
conditional form is proven~\cite{ouroboros25}, and the weaker-condition
uniqueness question remains open. Byzantine safety of the Khipu consensus is
\textbf{Conjecture~2}, open. Neither open question is used in the present
results.

\section{The honesty doctrine and proof tiers}\label{sec:doctrine}

Every claim in this paper carries exactly one epistemic tier.

\begin{itemize}
\item \textbf{[locked / kernel-verified]} --- proven in Lean~4 with only the
  trusted core axioms $\kernelaxioms$; \texttt{sorry}-free; pinned; part of the
  \emph{locked} count. Locked-proven $=$ \textbf{exactly eight}.
\item \textbf{[experimental, axiom-free, CI-green]} --- \texttt{sorry}-free with
  only core axioms, green in continuous integration (CI), residing in the
  \emph{experimental} tier; \textbf{never folded into the locked eight}.
\item \textbf{[axiom-gated]} --- \texttt{sorry}-free \emph{given} an explicitly
  declared idealizing axiom (e.g.\ hash collision-resistance), with that axiom
  in the \texttt{\#print axioms} ledger.
\item \textbf{[Conjecture]} --- a research hypothesis, never called a theorem.
  Conjecture~1 is disproved as originally stated, with a weaker-condition
  uniqueness question still open; Khipu safety is Conjecture~2.
\end{itemize}

\noindent The historical Doctrine-v11 count baseline is pinned at commit
\texttt{c7c0ba17} ($749$ declarations $/\,14$ unique axiom names $/\,163$
tracked \texttt{sorry} obligations). The claim-bearing F4/F7/F22 proof terms
and the eight-entry disclosure registry used by this paper are pinned at the
later proof snapshot \texttt{5cfaf9a3b55fd75e8a2a1197502a93dd8ca68be3}.
They are imported into the build and checked by \texttt{lake build}; the other
$\sim$185 CI-green declarations remain explicitly experimental. This tier
discipline --- a small locked registry and a large, honestly-labelled
experimental frontier --- follows the spirit of the seL4
machine-checked-systems tradition~\cite{sel4}, in which a precisely bounded
trusted base is the unit of trust.

\section{The receipt substrate}\label{sec:substrate}

We model the governance record as three cooperating structures.

\begin{definition}[Receipt DAG]
A \emph{receipt} is a node $0,1,\dots,n-1$ in emission order. A reference from a
later receipt to an earlier one is a directed edge $(\mathrm{src},
\mathrm{dst})$. The edge set is $\texttt{KhipuEdges} := \texttt{List}\,(\mathbb{N}
\times \mathbb{N})$. The append-only construction enforces the
\emph{backward-edge invariant} $\mathrm{dst} < \mathrm{src}$: a receipt may
reference only strictly-earlier receipts.
\end{definition}

\begin{definition}[Chaski channel]
Inter-organ messages flow over a FIFO queue $\texttt{ChaskiQueue} :=
\texttt{List}\,\mathbb{N}$. \emph{Enqueue} appends a message to the back (send);
\emph{dequeue}/\emph{drain} serves from the front (receive).
\end{definition}

\begin{definition}[Emit log]
The emit log after $n$ emissions is $\texttt{seqLog}\,n := \texttt{List.range}\,n
= [0,1,\dots,n-1]$; emitting appends the next sequence number.
\end{definition}

These three structures correspond, respectively, to the provenance graph, the
nervous/messaging plane, and the append-only ledger of the SZL ``living
anatomy'' substrate. We now prove their governing invariants.

\section{System boundary, threat model, and proof obligations}
\label{sec:threat}

The formal results are small enough to state exactly, but their operational
meaning depends on a boundary that is often left implicit. We therefore define
the principals, the attack surface, and the refinement obligations connecting
the Lean objects to a deployed receipt service before presenting the proofs.
This section is part of the claim: it prevents a theorem about lists from being
misread as an end-to-end theorem about a networked system.

\subsection{Principals and execution boundary}

An execution involves four logical principals. The \emph{emitter} $E$ is an
agent or tool runner that proposes an action and its evidence. The
\emph{recorder} $S$ assigns a sequence number, stores the receipt, and links it
to prior receipts. The \emph{channel} $C$ carries inter-organ notifications.
The \emph{auditor} $A$ reconstructs the history from the stored receipts and
messages. These are logical roles rather than necessarily distinct machines.
In a minimal deployment $E$ and $S$ may share a process; in a stronger one, $S$
is an independently attested service outside the agent boundary.

The Lean development begins \emph{after} an input has crossed into three typed
objects: an edge list, a queue, and a sequence log. The theorems do not parse
JSON, authenticate a caller, select a database transaction isolation level, or
prove that a network transport implements a queue. Those are refinement
obligations. The verified core establishes what follows if the adapter exposes
objects satisfying the stated constructors and invariants.

\subsection{Adversary model}

We consider an adversary that may control an agent prompt, tool output, receipt
payload, message contents, timing, and attempted references to other receipts.
It may try to create circular justification, reorder notifications, reuse a
sequence number, omit a message, or present a stale snapshot as current. The
adversary may also read all public proof artifacts. The guarantees are safety
properties and do not rely on obscurity.

The adversary is \emph{not} assumed able to alter a kernel-checked proof while
retaining the same Git commit identifier, subvert the Lean kernel, or violate
the arithmetic ordering of natural numbers. Cryptographic collision resistance
is deliberately outside the structural model and is isolated in
\S\ref{sec:binding}. Availability, confidentiality, fair scheduling, Byzantine
agreement, wall-clock ordering, and faithful execution of an unverified adapter
are also out of scope. This exclusion is not a weakness hidden in footnotes: it
determines the exact assurance supplied by F4, F7, and F22.

\subsection{Safety obligations}

Let $H=(R,E_H,Q,L)$ be a recorded history, where $R$ is the finite set of
receipt identifiers, $E_H\subseteq R\times R$ is the reference relation, $Q$
is the channel state, and $L$ is the emission log. The verified model discharges
the following obligations.

\begin{description}
\item[O1 -- backward references.] Every stored edge $(s,d)$ satisfies $d<s$.
This is the single local admission condition from which global acyclicity is
derived.
\item[O2 -- fresh-node append.] A batch appended for fresh identifier $k$ has
source $k$ on every new edge and targets only identifiers smaller than $k$.
\item[O3 -- FIFO construction.] Send appends to the back of the queue and
receive removes from the front; a batch preserves list order.
\item[O4 -- sequence construction.] The $n$-entry log is
$[0,\ldots,n-1]$ and the next emit appends $n$.
\end{description}

The model then yields the global conclusions in
Table~\ref{tab:assumption-claim}. Notice that no conclusion silently exceeds
its premise.

\begin{table}[t]
\centering
\small
\begin{tabular}{@{}p{0.20\linewidth}p{0.29\linewidth}p{0.39\linewidth}@{}}
\toprule
Result & Consumed obligation & Certified conclusion \\
\midrule
F4 & O1; O2 for append & No non-empty directed path returns to its source;
append preserves that fact. \\
F7 & O3 & Draining a batch sent to an empty channel returns the same ordered
list, including option-valued positional equality. \\
F22 & O4 & A new emission extends the prior log and values strictly increase
with position. \\
Count theorem & Published disclosure list & The registry contains eight
entries and every disclosed axiom name is on the allowed core list. \\
\bottomrule
\end{tabular}
\caption{Assumption-to-claim map. The count theorem concerns the disclosed
registry; per-declaration \texttt{\#print axioms} output remains the ground
truth for kernel dependencies.}
\label{tab:assumption-claim}
\end{table}

\subsection{A compositional assurance statement}

The three proofs compose without depending on one another. If a recorder
enforces O1--O4 at its trusted boundary, then every accepted history has an
acyclic reference relation, its modeled notification batch is observed in send
order, and its modeled emission log is a strict prefix-preserving sequence.
This is a conjunction of safety properties, not a new consensus theorem. In
particular, the conclusion does not say that every proposed receipt is accepted,
that every accepted receipt is semantically truthful, or that two Byzantine
replicas select the same history. It says that any history admitted through the
specified constructors has the three structural properties. This is the useful
contract for a refinement proof: an implementation can target a small, explicit
interface instead of an informal phrase such as ``tamper-proof audit trail.''

\section{F4 --- Acyclicity of the receipt DAG, preserved under append}\label{sec:f4}

\subsection{Reachability and the key lemma}

\begin{definition}[Reachability]
$\texttt{KhipuStep}\;es\;a\;b := (a,b)\in es$ is one directed edge.
$\texttt{KhipuReach}\;es$ is its non-reflexive transitive closure: a directed
path of length $\geq 1$.
\end{definition}

\begin{lemma}[Paths strictly decrease the index]\label{lem:decrease}
\emph{[locked]} If $es$ satisfies the backward-edge invariant
$\forall e \in es,\ e_2 < e_1$, then $\texttt{KhipuReach}\;es\;a\;b
\Rightarrow b < a$.
\end{lemma}

\begin{proof}
By induction on the reachability derivation. In the \texttt{base} case a single
edge $(a,b)$ gives $b < a$ directly from the invariant. In the \texttt{trans}
case, an edge $a\to b$ and a path $b\rightsquigarrow c$ give, by the inductive
hypothesis $c < b$ and the invariant $b < a$, the conclusion $c < a$ via
transitivity of $<$. \qed
\end{proof}

\begin{theorem}[Acyclicity, F4]\label{thm:f4-acyclic}
\emph{[locked, \texttt{f4\_khipu\_no\_cycle}]} Under the backward-edge invariant,
the transitive closure is irreflexive: $\forall a,\ \neg\,
\texttt{KhipuReach}\;es\;a\;a$. The receipt DAG has no directed cycle.
\end{theorem}

\begin{proof}
A self-reachability $\texttt{KhipuReach}\;es\;a\;a$ would, by
Lemma~\ref{lem:decrease}, give $a < a$, contradicting irreflexivity of $<$
(\texttt{Nat.lt\_irrefl}). \qed
\end{proof}

\begin{theorem}[Append preserves acyclicity, F4 main]\label{thm:f4-append}
\emph{[locked, \nolinkurl{f4_khipu_dag_acyclic_preserved}]} Let $es$ satisfy the
invariant. Appending a fresh maximal node $k$ together with edges whose source is
$k$ and whose targets are all $< k$ yields an edge set $es \mathbin{+\!\!+}
\mathit{newEdges}$ that still has no directed cycle.
\end{theorem}

\begin{proof}
First, $es \mathbin{+\!\!+} \mathit{newEdges}$ still satisfies the backward-edge
invariant (\nolinkurl{f4_khipu_append_preserves_invariant}): for an old edge the
invariant holds by hypothesis; for a new edge $e$ we have $e_1 = k$ and
$e_2 < k$, hence $e_2 < e_1$. Acyclicity then follows from
Theorem~\ref{thm:f4-acyclic} applied to the extended set. \qed
\end{proof}

\subsection{Why this is non-vacuous}

A prior formalization stated F4 as
$\texttt{f4\_khipu\_dag\_acyclic}\,(k\,t)\,(h:t<k) : t<k \wedge k\neq t$, which
merely \emph{repackages its hypothesis} and proves nothing global. An internal
audit (\S\ref{sec:nonvacuity}) flagged it as vacuous. The statements above
instead quantify over an \emph{arbitrary} edge set and an \emph{arbitrary}
reachability derivation, and conclude a genuine global property ---
irreflexivity of the transitive closure --- which cannot be obtained without the
inductive argument of Lemma~\ref{lem:decrease}. This distinction is the heart of
honest formalization: a theorem must not be derivable from its own hypothesis by
unfolding.

\section{F7 --- FIFO reception ordering of the Chaski channel}\label{sec:f7}

\begin{definition}[Channel operations]
$\texttt{chaskiEnqueue}\;q\;m := q \mathbin{+\!\!+} [m]$ (send to back);
$\texttt{chaskiEnqueueAll}\;q\;xs := q \mathbin{+\!\!+} xs$ (batch send);
$\texttt{chaskiDrain}$ serves the front recursively, collecting reception order.
\end{definition}

\begin{lemma}[Drain preserves order]\label{lem:drain}
\emph{[locked, \texttt{f7\_chaski\_drain\_eq}]} $\texttt{chaskiDrain}\;q = q$.
\end{lemma}

\begin{proof}
Structural induction on $q$: the empty queue drains to the empty list; for
$m :: ms$, draining gives $m :: \texttt{chaskiDrain}\;ms = m :: ms$ by the
inductive hypothesis. \qed
\end{proof}

\begin{theorem}[Reception order equals send order, F7 main]\label{thm:f7}
\emph{[locked, \texttt{f7\_chaski\_fifo\_order}]} Enqueuing a batch
$\mathit{msgs}$ onto an empty channel and then draining yields exactly
$\mathit{msgs}$: $\texttt{chaskiDrain}\,(\texttt{chaskiEnqueueAll}\,[]\,
\mathit{msgs}) = \mathit{msgs}$.
\end{theorem}

\begin{proof}
$\texttt{chaskiEnqueueAll}\,[]\,\mathit{msgs} = []\mathbin{+\!\!+}\mathit{msgs}
= \mathit{msgs}$, after which Lemma~\ref{lem:drain} gives
$\texttt{chaskiDrain}\,\mathit{msgs} = \mathit{msgs}$. \qed
\end{proof}

\begin{corollary}[Positional FIFO witness]\label{cor:f7pos}
\emph{[locked, \nolinkurl{f7_chaski_fifo_positional}]} For every index $i$, the
$i$-th received message equals the $i$-th sent message:
\[
(\texttt{chaskiDrain}\,(\texttt{chaskiEnqueueAll}\,[]\,\mathit{msgs}))[i]?
= \mathit{msgs}[i]?
\]
(option-valued indexing, so no in-range side goal arises).
\end{corollary}

A prior formalization stated F7 as the tautology $\mathit{msgs} = \mathit{msgs}$.
The statement above instead equates the output of a non-trivial drain over an
enqueued batch with the input, which is a genuine order-preservation claim;
Corollary~\ref{cor:f7pos} makes it an element-by-element witness.

\section{F22 --- Append-only emit monotonicity}\label{sec:f22}

\begin{theorem}[Emit appends the next sequence number, F22]\label{thm:f22-append}
\emph{[locked, \nolinkurl{f22_emit_appends_length}]} $\texttt{seqLog}\,(n+1) =
\texttt{seqLog}\,n \mathbin{+\!\!+} [n]$.
\end{theorem}

\begin{theorem}[Strict positional monotonicity, F22 main]\label{thm:f22}
\emph{[locked, \texttt{f22\_khipu\_emit\_monotone}]} For $i < j < n$, the entry at
position $i$ is strictly less than the entry at position $j$ in
$\texttt{seqLog}\,n$.
\end{theorem}

\begin{proof}
$\texttt{seqLog}\,n = \texttt{List.range}\,n$, so its element at position $p$ is
$p$ itself (\texttt{List.getElem\_range}); hence the entry comparison reduces to
$i < j$. \qed
\end{proof}

For this constructed list, strict positional monotonicity and the
\texttt{range\_succ} prefix equation supply the modeled append-only property:
the next state preserves the prior prefix and adds one fresh sequence value.
This does not by itself prevent a storage implementation from overwriting
history. Together with F4 and F7, it supplies a totally ordered, non-circular
model that a deployed governance record must refine.

\section{Derived consequences and an admission algorithm}
\label{sec:derived}

The named Lean declarations are intentionally minimal. Several useful
consequences follow on paper from the same definitions; we record them here so
that system designers know which claims are immediate and which require new
formal work. Unless a Lean identifier is given, a consequence in this section
is a mathematical corollary, not an additional counted theorem.

\subsection{A canonical topological order}

Under the backward-edge invariant, every path strictly decreases the natural
identifier. Therefore the descending identifier order is a topological order
for the edge orientation used here: every source appears before every target.
Equivalently, ascending order is a valid construction order because a node may
refer only to nodes already constructed. This gives an auditor a simple replay
schedule without running a general topological-sort algorithm. It also gives a
termination measure for recursive ancestry traversal: the current identifier
strictly decreases on every recursive call.

The convenience has a cost. Identifier order becomes semantically relevant to
the construction, so importing an arbitrarily labelled DAG requires a trusted
topological relabelling step. The present artifact neither implements nor
verifies such an importer. A system that accepts external graph identifiers
should preserve an internal monotone index and treat the external identifier as
payload, rather than comparing opaque strings as though they were the verified
order.

\subsection{Prefix and uniqueness consequences of the range log}

For $m\leq n$, \texttt{seqLog}$\,m$ is a prefix of
\texttt{seqLog}$\,n$. This follows by repeated application of
Theorem~\ref{thm:f22-append}. Every value in \texttt{seqLog}$\,n$ is smaller
than $n$, every position $i<n$ contains $i$, and no value appears twice. Thus
the length, next identifier, and set of emitted identifiers determine one
another in the model. These consequences justify an implementation invariant
in which the transactionally stored next counter equals the committed receipt
count, but only when deletion is impossible or separately represented by a
tombstone. Physical deletion would break the abstraction even if the remaining
rows were sorted.

\subsection{Observational content of FIFO equality}

List equality in F7 is stronger than equality of message multisets. It implies
equal length and equal option-valued lookup at every index, so order, loss, and
duplication are all observable within the modeled batch. Conversely,
positional equality for all natural indices implies list equality because both
out-of-range lookups return \texttt{none}. The theorem therefore supports two
styles of runtime test: compare whole serialized batches, or compare length and
each indexed message. The latter produces a more diagnostic failure receipt.

The equality is deliberately parameterized by message identifiers of type
\texttt{Nat}. It does not inspect payload equality, cryptographic identity, or
delivery side effects. A realistic adapter should map each durable message to a
unique natural identifier and separately attest the identifier-to-payload
binding.

\subsection{Reference admission procedure}

The following procedure is the executable shape suggested by F4 and F22. It is
pseudocode, not verified application code.

\begin{lstlisting}[language=]
appendReceipt(payload, targets):
  begin serializable transaction
    k := allocateNextSequence()
    if any target is not a committed integer in [0, k):
      rollback and return REJECTED_INVALID_REFERENCE
    store receipt(k, payload)
    for target in deduplicate(targets):
      store edge(k, target)
    append emit_log(k)
  commit
  return ACCEPTED(k)
\end{lstlisting}

If the transaction begins from a state refining the Lean model and allocation
is fresh, the local range check establishes O2 for all new edges, the append to
the emit log establishes O4, and Theorems~\ref{thm:f4-append} and
\ref{thm:f22-append} preserve the two global invariants. Deduplication is not
needed for acyclicity but avoids ambiguous edge multiplicity because the Lean
representation is a list while the paper often speaks set-theoretically.

With an indexed receipt table, checking $r$ targets takes $O(r)$ membership and
order checks; appending the rows is also $O(r)$. No traversal of the existing
graph is required because the well-founded index is a local certificate of
global acyclicity. This is an algorithmic advantage of the construction, not a
measured performance claim. Storage engines, transaction contention, and
signature generation determine actual latency and are not benchmarked here.

\subsection{Composition with the channel}

After commit, the service may enqueue a notification containing $k$. If this
enqueue and database commit are not atomic, a transactional outbox is a common
refinement: commit an outbox row in the receipt transaction, publish rows in
increasing outbox position, and record consumption offsets. F7 can then model
one ordered outbox partition. Without such a bridge, an acyclic and monotone
database can coexist with lost or reordered notifications. The proof family
therefore recommends an architecture but does not certify a particular broker.

\subsection{Worked trace}

Consider six receipts admitted in ascending sequence order. Receipt zero is a
root. Receipt one cites zero; receipt two cites zero; receipt three cites one
and two; receipt four cites two; and receipt five cites three and four. The edge
list is
\[
[(1,0),(2,0),(3,1),(3,2),(4,2),(5,3),(5,4)].
\]
Every target is smaller than its source. A path from five may visit
$5\to3\to1\to0$ or $5\to4\to2\to0$, but the identifier strictly decreases at
each step, so neither path can return to five. The emission log is
$[0,1,2,3,4,5]$ and the notification batch for the last three admissions is
$[3,4,5]$.

\begin{table}[h]
\centering
\small
\begin{tabular}{@{}ccp{0.42\linewidth}p{0.28\linewidth}@{}}
\toprule
$k$ & Targets & Local admission result & Global consequence \\
\midrule
0 & $\varnothing$ & Accept root & Empty graph is acyclic. \\
1 & $\{0\}$ & $0<1$ & Append preserves F4. \\
2 & $\{0\}$ & $0<2$ & Independent branch is allowed. \\
3 & $\{1,2\}$ & Both targets $<3$ & Merge node remains acyclic. \\
4 & $\{2\}$ & $2<4$ & Append and log prefix preserved. \\
5 & $\{3,4\}$ & Both targets $<5$ & Two ancestry paths terminate. \\
5 (retry) & $\{5\}$ & Reject: $5\not<5$ & No self-loop; transaction rolls back. \\
6 & $\{7\}$ & Reject: $7\not<6$ & Future reference cannot enter history. \\
\bottomrule
\end{tabular}
\caption{A concrete admission trace. The last two rows are attempted writes,
not committed receipts; a correct adapter does not consume their sequence
numbers unless its external contract explicitly records rejected attempts.}
\label{tab:trace}
\end{table}

Suppose the channel sends notifications $3$, $4$, and $5$ by applying
append-to-back three times. Starting from the empty list, its successive states
are $[]$, $[3]$, $[3,4]$, and $[3,4,5]$. Front-drain produces $3$, then $4$,
then $5$. F7 identifies the whole received list with the sent list; its
positional corollary separately witnesses equality at indices zero, one, two,
and every out-of-range index.

Now mutate the final accepted edge to $(5,6)$. The new edge does not itself
create a cycle in this particular finite graph because six is absent, but it
violates the constructive invariant and must still be rejected. This illustrates
the difference between checking the \emph{current} graph for a cycle and
enforcing an inductive admission rule. The local rule deliberately rejects some
states that happen to be acyclic so that every future append can be certified
without traversing prior history. The design trades permissiveness for a small,
stable proof obligation and constant-history admission checks.

The trace also exposes a policy choice not fixed by F22: whether rejected
attempts receive identifiers. The verified \texttt{seqLog} models committed
emissions only. If an organization needs a forensic log of rejected attempts,
that log should use a separate monotone namespace or represent acceptance as a
status field without deleting identifiers. Conflating attempted and committed
sequences would make the paper's abstraction ambiguous.

\section{The locked count is exactly eight --- itself a theorem}\label{sec:count}

We render the count discipline as a machine-checked statement. The disclosed
axiom set of each locked theorem is modelled as a concrete list of axiom names;
the kernel base is $\texttt{leanKernelAxioms} = \kernelaxioms$.

\begin{lstlisting}[language=]
def lockedDisclosed : List (String x List String) :=
  [ ("F1",  ["propext","Classical.choice","Quot.sound"]),
    ("F4",  ["propext","funext","Classical.choice","Quot.sound"]),
    ("F7",  ["propext","funext","Classical.choice","Quot.sound"]),
    ("F11", ["propext","Classical.choice","Quot.sound"]),
    ("F12", ["propext","Classical.choice","Quot.sound"]),
    ("F18", ["propext","funext","Classical.choice","Quot.sound"]),
    ("F19", ["propext","Classical.choice","Quot.sound"]),
    ("F22", ["propext","funext","Classical.choice","Quot.sound"]) ]

theorem locked_axiom_sets_kernel_only :
    lockedDisclosed.all (fun p => axiomsAllowed p.2) = true := by decide

theorem locked_count_eight : lockedDisclosed.length = 8 := by decide
\end{lstlisting}

\noindent\texttt{locked\_axiom\_sets\_kernel\_only} certifies, by \texttt{decide},
that every \emph{published disclosure-name list} lies inside the allowed core
name set. It does not inspect theorem dependencies. \texttt{locked\_count\_eight}
certifies that the published registry contains exactly eight entries. Five
\emph{independent} count theorems
(in \texttt{Wave8/Wave9/Wave10/Wave11/AxiomDisclosure} and
\texttt{Uniqueness/AxiomCheck}) were moved $5 \to 8$ in lockstep so that the
served count cannot silently diverge from the proven count. This is the count
becoming an object the kernel checks, rather than a number in prose.

\begin{remark}[A meta-level caveat, disclosed]
\texttt{locked\_axiom\_sets\_kernel\_only} is a property of the disclosed axiom
\emph{names}; it does not, and cannot by itself, re-verify the underlying kernel
dependencies. A release audit must capture per-declaration
\texttt{\#print axioms} output from the named snapshot and reject
\texttt{sorryAx} or undeclared project axioms. We state both so the reader can
distinguish the registry meta-claim from the object-level proof audit.
\end{remark}

\section{Cryptographic binding is axiom-gated (not a hardness proof)}\label{sec:binding}

The three invariants above are about \emph{structure}, not \emph{secrecy}. In
deployment, each receipt also carries a SHA-256 hash-chain entry with a
per-batch Merkle root, and F18 (Reed--Solomon coding, locked) governs erasure
tolerance. Hash-chain \emph{binding} --- that you cannot alter a past entry
without breaking the root --- is proven only \textbf{[axiom-gated]} under a
declared \texttt{sha256\_collision\_resistant} axiom plus Merkle
collision-resistance / domain-separation idealizations, all disclosed in the
\texttt{\#print axioms} ledger. We do not prove cryptographic hardness; we
isolate it as an explicit, auditable assumption. The structural theorems F4/F7/F22
depend on \emph{none} of these idealizations.

\section{The non-vacuity audit as a methodology}\label{sec:nonvacuity}

Both F4 and F7 had previously-landed ``proofs'' that were honest in tone but
mathematically vacuous: F4 repackaged its hypothesis; F7 was $\mathit{msgs} =
\mathit{msgs}$. An internal audit identified these and the count was \emph{not}
advanced to eight until genuine statements replaced them on 2026-06-10. We
extract a reusable test: a candidate theorem is suspect when its conclusion is
syntactically derivable from its hypotheses by unfolding definitions, or when it
has the form $x = x$ after reduction. The repair is to (i) introduce a data model
rich enough to express the intended global property (here, an edge-set DAG and a
FIFO queue), (ii) quantify over arbitrary objects and arbitrary derivations, and
(iii) require the proof to consume genuine structure (induction on a reachability
derivation; structural induction on a queue). We report this because vacuous
formalizations are a recurring failure mode in applied verification, and
detecting them is part of an honest pipeline.

\section{Countermodels and falsification-oriented review}
\label{sec:countermodels}

A proof artifact is easier to trust when the reader can also see what would
break the claim. We therefore accompany each positive theorem with a small
countermodel or mutation that its hypotheses reject. These are not empirical
benchmarks and are not presented as evidence about throughput. They are
semantic negative controls: each identifies a nearby, plausible statement that
is false.

\subsection{F4 negative controls}

The backward-edge invariant is sufficient but not decorative. Consider the two
edges $(1,0)$ and $(0,1)$. The first is a valid backward reference; the second
is forward and violates $1<0$. Together they form the cycle
$1\rightarrow0\rightarrow1$. If the admission adapter checks only that a target
exists, rather than that it is strictly earlier, F4 does not apply. Likewise,
allowing a fresh node $k$ to emit an edge with target $k$ introduces an immediate
self-loop. These examples show why the append theorem needs both clauses of
its hypothesis: all new sources equal the fresh node and all new targets are
strictly smaller.

The converse does not hold: a DAG can contain an edge whose numeric target is
larger than its source under an arbitrary labelling. The backward-index
invariant is therefore a constructive discipline, not a characterization of
all DAGs. A deployment may topologically relabel an existing DAG before import,
but the present theorem does not prove that relabelling algorithm. This
distinction matters when migrating a legacy audit graph.

\subsection{F7 negative controls}

Three mutations expose the scope of the FIFO theorem. First, define
$\texttt{enqueueFront}\;q\;m := m::q$ while keeping front-drain; two sends
$a,b$ are then received $b,a$. Second, retain append-to-back but define drain to
reverse the queue; order again fails. Third, keep the correct operations but
allow a transport to acknowledge a message before durably enqueueing it; the
list theorem remains true of the in-memory list while the system-level claim
can fail after a crash. The first two are rejected by O3; the third is a
refinement failure outside the Lean model.

F7 proves order preservation, not exactly-once delivery. The equation for a
single batch rules out loss and duplication \emph{within that modeled drain},
but it does not assign globally unique message identifiers, persist offsets, or
deduplicate retries across drains. An implementation that needs exactly-once
effects must add idempotency and durable-offset obligations rather than
renaming FIFO as exactly-once.

\subsection{F22 negative controls}

For $L=[0,2,1]$, all entries are distinct, yet strict positional monotonicity
fails between positions one and two. For $L=[0,1,1]$, append has occurred but
freshness fails. For $L'=[0,1,2]$ obtained by replacing an earlier log
$[0,9]$, the final list is monotone even though the history was overwritten.
F22 avoids the last ambiguity because the construction is not an arbitrary
monotone list: \texttt{seqLog} is definitionally \texttt{List.range}, and the
append theorem connects state $n$ to state $n+1$. In a database refinement, the
analogous obligation is transactional prefix preservation, not merely a query
that returns sorted rows.

\subsection{A non-vacuity rubric}

The audit used the following review rubric. It is intentionally mechanical so
that it can be automated in part.

\begin{enumerate}
\item Normalize the statement enough to detect conclusions of the form $x=x$,
$P\to P$, or a conjunction containing a hypothesis verbatim.
\item Remove each hypothesis in turn and ask whether the conclusion still
follows by reflexivity, simplification, or constructor application.
\item Construct the smallest mutation of the data model that should violate the
intended property; confirm that at least one explicit premise excludes it.
\item Check that the proof consumes the relevant structure. Induction over a
path should inspect the path; a FIFO proof should unfold or induct over the
queue; an append theorem should relate pre- and post-state.
\item Separate theorem non-vacuity from model adequacy. A genuine theorem over
an inadequate model is still inadequate for deployment.
\end{enumerate}

The rubric would have rejected the earlier F4 and F7 statements before merge.
It also catches a subtler failure mode: a theorem can mention a graph or queue
type while never using its edges or operations. This ``phantom structure'' test
is now part of the recommended review procedure for future formula promotions.

\section{Refinement from Lean objects to a receipt service}
\label{sec:refinement}

The verified substrate becomes operational only through a refinement map
$\rho$ from concrete state to the three Lean objects. Let a concrete database
state contain receipt rows, reference rows, message rows, and an emission
counter. The abstraction map sorts or indexes these rows into
$(\texttt{KhipuEdges},\texttt{ChaskiQueue},\texttt{seqLog})$. Five conditions
are sufficient to make the structural theorems relevant to that service.

\begin{description}
\item[R1 -- atomic admission.] Allocating $k$, validating its references, and
committing the receipt and edges occur in one transaction. Otherwise a crash
can expose edges without a receipt or reuse $k$.
\item[R2 -- canonical identifiers.] The database comparison used by the
adapter agrees with the natural-number order used in Lean. Lexicographic string
ordering is not a valid substitute (for example, ``10'' precedes ``2'').
\item[R3 -- durable FIFO state.] Queue order is represented by a durable,
unique position or broker offset and dequeue commits consumption atomically.
Wall-clock timestamps alone are insufficient because ties and clock skew do
not define FIFO.
\item[R4 -- prefix-preserving reads.] An audit snapshot corresponds to a
committed prefix. Pagination, replicas, and caches must not splice rows from
different snapshots into a synthetic history.
\item[R5 -- attested adapter.] The binary or image implementing $\rho$ is
identified by a build provenance statement and the deployment selects that
identified artifact. A proof about one adapter does not cover an untracked
replacement.
\end{description}

These conditions can be tested without trusting an agent's prose. R1 and R4
admit transactional integration tests with injected failures. R2 admits
property tests over boundary identifiers. R3 admits enqueue/dequeue histories
with retries. R5 admits digest comparison against a signed build attestation.
None of these tests replaces the Lean proofs; conversely, the Lean proofs do
not replace these tests. The useful assurance case is heterogeneous: deductive
proof for mathematical structure, adversarial tests for the adapter, and
attestation for artifact identity.

\subsection{Crash and concurrency cases}

Concurrency introduces two important cases. If two emitters both observe
counter $n$ and allocate it without serialization, F22's concrete premise is
already false. A unique constraint may reject one write, but the adapter must
then retry from the new counter rather than silently dropping the receipt. If a
message is delivered to two consumers, FIFO per consumer does not establish a
single global receive order. A deployment must state whether $Q$ denotes one
partition, one consumer group, or a merged stream; the refinement proof is per
chosen ordering domain.

Crash recovery must preserve a prefix. A write-ahead log or transactional
database may satisfy this naturally, while an eventually consistent object
store may require an explicit manifest whose publication is the commit point.
The theorem is agnostic about the mechanism. What matters is that recovery
cannot expose a state that maps to $[0,\ldots,n-2,n]$ or to a new edge without
its fresh source node.

\subsection{What an auditor should request}

For a concrete service, an auditor should request: the abstraction-map version;
the exact source commit and build digest; the database schema or broker ordering
contract; admission tests for self, forward, and missing references; a
concurrent-allocation test; a crash-recovery trace; the per-theorem
\texttt{\#print axioms} output; and the policy that keeps experimental
declarations outside the locked registry. This evidence is deliberately more
specific than a badge saying ``formally verified.''

\section{Artifact audit and reproducibility protocol}
\label{sec:artifact}

We performed a source-and-CI audit on 2026-07-15. It revealed an important
distinction that an earlier draft blurred: \texttt{c7c0ba17} is the frozen
Doctrine-v11 \emph{baseline} used for the canonical $749/14/163$ declaration,
axiom-name, and tracked-obligation counts, but it predates the non-vacuous
F4/F7/F22 artifact. The paper's proof snapshot is
\texttt{5cfaf9a3b55fd75e8a2a1197502a93dd8ca68be3} (2026-06-10).
The first proof implementation merged through pull request~185 at
\texttt{e6de491dc1eddf9d1a6ea1306a10c96e3971d888}; the later snapshot replaces
the vacuous statements, synchronizes the eight-entry registry, and is the
minimum commit that supports the claims made here. The baseline and proof
snapshot serve different purposes and must not be substituted for one another.

\begin{table}[t]
\centering
\small
\begin{tabular}{@{}p{0.18\linewidth}p{0.25\linewidth}p{0.45\linewidth}@{}}
\toprule
Artifact & Identifier & Role \\
\midrule
Doctrine baseline & \texttt{c7c0ba17c2eaec60} & Frozen canonical-count
reference; does not contain the final non-vacuous proof snapshot. \\
Initial proof merge & \texttt{e6de491dc1eddf9} & PR~185; CI \texttt{lake
build + numbers}, standalone build, doctrine, DCO, security, and DOI gates
reported success (the title-lint job was non-required and failed). \\
Claim snapshot & \texttt{5cfaf9a3b55fd75e} & Non-vacuous F4/F7, F22, and
locked-eight registry examined by this paper. \\
Paper repository & \nolinkurl{szl-holdings/szl-papers} & LaTeX, metadata, and
submission checklist; no arXiv identifier is claimed before moderation. \\
\bottomrule
\end{tabular}
\caption{Reproducibility identifiers and their non-interchangeable roles. Full
40-hex commit identifiers appear in the text and verification block.}
\label{tab:artifacts}
\end{table}

\subsection{Four verification levels}

\paragraph{Level 1: source presence.} At the claim snapshot, inspect
\nolinkurl{Lutar/Puriq/Formulas/ProvedFormulas.lean} and
\nolinkurl{Lutar/Wave8/AxiomDisclosure.lean}. Confirm that the theorem names and
statements match this paper and that the file contains no \texttt{sorry}.

\paragraph{Level 2: focused kernel check.} Run Lean on the self-contained
proved-formulas module. This is faster and narrows failures to the artifact
under review. It checks elaboration and proof terms but is not a substitute for
the repository build.

\paragraph{Level 3: repository build.} Run \texttt{lake build} with the pinned
toolchain. This checks the import graph used by CI. PR~185's public workflow
record reports both \texttt{lake build + numbers} and the separate build job as
successful at the merge commit.

\paragraph{Level 4: dependency disclosure and mutation.} Capture
\texttt{\#print axioms} for every promoted declaration, search the output for
\texttt{sorryAx} and project axioms, and run negative mutations: reverse an
edge, enqueue at the front, or replace \texttt{List.range} with a duplicate.
The mutations should make the relevant theorem unprovable or its premise
unsatisfied. A green build without this scope review can still certify the
wrong statement.

\subsection{Reproducibility commands}

The following sequence separates the baseline check from the claim check.

\begin{lstlisting}[language=]
git clone https://github.com/szl-holdings/lutar-lean.git
cd lutar-lean

# The exact artifact supporting F4/F7/F22 and the eight-entry registry.
git checkout 5cfaf9a3b55fd75e8a2a1197502a93dd8ca68be3
lake build

# Inspect the source declarations and disclosure registry.
grep -n "theorem f4_khipu\|theorem f7_chaski\|theorem f22_" \
  Lutar/Puriq/Formulas/ProvedFormulas.lean
grep -n "locked_count_eight\|locked_axiom_sets_kernel_only" \
  Lutar/Wave8/AxiomDisclosure.lean

# The frozen canonical-count baseline is a separate historical check.
git checkout c7c0ba17c2eaec60ad38ea9172b4a0d9ca0b582f
python .github/scripts/lean_numbers.py
\end{lstlisting}

Because the focused proof module and the count registry are distinct files, a
reviewer should not infer the former's actual axiom dependencies solely from
the latter's list of disclosed strings. The per-declaration kernel output is
the authoritative dependency record; the list theorem is a machine-checked
consistency control over published metadata.

\section{Assurance ledger and residual risk}
\label{sec:ledger}

Table~\ref{tab:ledger} summarizes what a release may and may not say after these
results. It is designed to be copied into a deployment review without turning a
structural theorem into a marketing superlative.

\begin{table}[h]
\centering
\small
\begin{tabular}{@{}p{0.22\linewidth}p{0.32\linewidth}p{0.34\linewidth}@{}}
\toprule
Control & Supported statement & Residual risk / prohibited inference \\
\midrule
Receipt references & Accepted modeled histories are acyclic under the
backward-edge admission rule. & Does not authenticate payloads, prove semantic
truth, or cover an adapter that fails to enforce the rule. \\
Channel order & A modeled batch drains in the same order it was enqueued. & No
cross-partition total order, durability, exactly-once effect, or fair delivery. \\
Emission order & \texttt{seqLog} extends by the next natural and is strictly
increasing by position. & No wall-clock synchronization, database isolation, or
automatic tamper evidence. \\
Axiom registry & Eight disclosed entries pass a deny-by-default name check. &
The string list is not a reimplementation of \texttt{\#print axioms}. \\
CI record & The merge workflow reports successful build, doctrine, DCO, and
security jobs for PR~185. & CI proves the checked commit and workflow, not an
unidentified deployment. \\
\bottomrule
\end{tabular}
\caption{Assurance ledger for release and audit language.}
\label{tab:ledger}
\end{table}

The highest residual risk is refinement drift: a deployed service can evolve
while the formal abstraction stays fixed. We recommend versioning the
abstraction map as part of the receipt schema and making its version a required
field in build provenance. A second risk is registry drift: a prose page can
say ``eight'' while the executable set differs. The five count theorems reduce
this risk within the source tree, but cross-repository publishing still needs a
link checker or signed manifest. A third risk is over-composition: acyclicity,
FIFO, and monotonicity together do not imply Byzantine agreement or
cryptographic immutability. Those claims remain separately gated.

\section{Related work}\label{sec:related}

\paragraph{Append-only ledgers and provenance.} Hash-chained, append-only logs
underpin transparency systems and certificate transparency; our backward-edge
DAG is the order-theoretic skeleton beneath such logs. The supply-chain
provenance layer that attests the \emph{artifact} producing receipts follows the
SLSA build track~\cite{slsa} with DSSE/Sigstore keyless signing~\cite{sigstore};
those attestations are captured outside the agent's process, the top rung of the
governance enforceability ladder~\cite{armo2025}.

\paragraph{Machine-checked systems and BFT.} The discipline of a small,
precisely-bounded trusted base is the seL4 tradition~\cite{sel4}. The consensus
layer that decides which receipts are canonical is governed by quorum
intersection in the $n \geq 3f+1$ regime~\cite{pbft,lamport1982}, with Byzantine
safety formalized in the Velisarios Coq development~\cite{velisarios}; in the SZL
substrate that safety is conditional (Conjecture~2 unconditionally).

\paragraph{Graph structure of learning systems (lineage we engage as prior
art).} The receipt DAG and the broader SZL graph substrate connect to a body of
work on \emph{graph structure in neural and learning systems}: position-aware
graph neural networks (P-GNN)~\cite{you2019pgnn} and the GraphGym/SNAP
design-space program~\cite{you2020graphgym}, the relational-graph view of neural
network structure (graph2nn)~\cite{you2020graph2nn}, the expressivity ceiling of
message passing at $1$-WL~\cite{xu2019gin}, scientific machine learning at
the Polymathic AI program~\cite{mccabe2023multiphysics}, and recurrent-depth /
reasoning-by-iteration~\cite{mcleish2024abacus}. We cite these as related and
prior art that informs our \emph{thinking} about graph-structured governance; we
do \emph{not} claim any of them as a source of the present results, which are
elementary order theory machine-checked for an audit setting.

\section{Limitations and honest posture}\label{sec:limits}

\begin{itemize}
\item F4/F7/F22 are \emph{structural} guarantees over idealized data models
  (edge-set DAG, FIFO queue, range log); they do not by themselves guarantee that
  a deployed implementation populates these structures faithfully --- that is the
  job of the attested provenance layer.
\item Cryptographic binding is axiom-gated (\S\ref{sec:binding}); we prove no
  hardness.
\item $\Lambda$ \textbf{Conjecture~1 is disproved as stated}; only its
  weaker-condition uniqueness question remains open. Khipu Byzantine safety is
  \textbf{Conjecture~2}, open; neither open question is used
  here~\cite{ouroboros25}.
\item The $\sim$185 experimental CI-green theorems (Waves~11--23) are a separate
  tier and are \emph{never} folded into the locked eight.
\item Trust is never reported as $100\%$. No fabricated benchmarks or citations.
\end{itemize}

\section{Conclusion}\label{sec:conclusion}

We have given a machine-checked, non-vacuous foundation for three structural
invariants of a modeled AI decision-receipt substrate --- acyclicity (F4), FIFO
batch reception (F7), and range-log monotonicity (F22) --- with proof terms
checked by Lean~4 at the named snapshot. The associated disclosure registry
contains exactly eight entries, with its count rendered as a theorem; its
allowed-name check is not a substitute for captured per-declaration
\texttt{\#print axioms} output. The mathematics is elementary by design: the
contribution is an honest, auditable governance receipt layer for
post-deterministic agents, a kernel-checked count discipline, and a reusable
methodology for detecting and repairing vacuous formalizations. The harder
guarantees --- weaker-condition aggregator uniqueness and Byzantine safety ---
remain explicitly open questions, and we say so.

\paragraph{Reproducibility.} The exact proof snapshot is
\texttt{5cfaf9a3b55fd75e8a2a1197502a93dd8ca68be3}; the distinct frozen
canonical-count baseline is \texttt{c7c0ba17c2eaec60ad38ea9172b4a0d9ca0b582f}.
Section~\ref{sec:artifact} gives the commands, source paths, CI record, and
reason the two commits must not be conflated. A valid release axiom audit must
capture per-declaration output reporting only the four allowed Lean core axioms
named in \S\ref{sec:doctrine}, no \texttt{sorryAx}, and no declared project
axiom for the promoted declarations.

\appendix

\section{Selected Lean artifact: acyclicity and append}
\label{app:f4}

The following excerpt is normalized only for line length; identifiers and proof
structure match
\texttt{Lutar/Puriq/Formulas/ProvedFormulas.lean} at the claim snapshot. It
shows why F4 is not a theorem about a Boolean ``DAG'' flag: reachability is an
inductive transitive closure and the proof consumes its derivation.

\begin{lstlisting}[language=]
abbrev KhipuEdges := List (Nat x Nat)

def KhipuBackwardInvariant (es : KhipuEdges) : Prop :=
  forall e in es, e.2 < e.1

def KhipuStep (es : KhipuEdges) (a b : Nat) : Prop :=
  (a, b) in es

inductive KhipuReach (es : KhipuEdges) : Nat -> Nat -> Prop
  | base  : KhipuStep es a b -> KhipuReach es a b
  | trans : KhipuStep es a b -> KhipuReach es b c ->
            KhipuReach es a c

theorem f4_khipu_reach_decreases
    (es : KhipuEdges) (inv : KhipuBackwardInvariant es)
    (r : KhipuReach es a b) : b < a := by
  induction r with
  | base h => exact inv _ h
  | trans h _ ih => exact Nat.lt_trans ih (inv _ h)

theorem f4_khipu_no_cycle
    (es : KhipuEdges) (inv : KhipuBackwardInvariant es)
    (a : Nat) : not KhipuReach es a a := by
  intro r
  exact Nat.lt_irrefl a (f4_khipu_reach_decreases es inv r)

theorem f4_khipu_append_preserves_invariant
    (es : KhipuEdges) (inv : KhipuBackwardInvariant es)
    (k : Nat) (newEdges : KhipuEdges)
    (hk : forall e in newEdges, e.1 = k and e.2 < k) :
    KhipuBackwardInvariant (es ++ newEdges) := by
  intro e he
  rcases List.mem_append.mp he with h | h
  . exact inv e h
  . have hk' := hk e h
    rw [hk'.1]
    exact hk'.2

theorem f4_khipu_dag_acyclic_preserved ... :
    forall a, not KhipuReach (es ++ newEdges) a a :=
  f4_khipu_no_cycle (es ++ newEdges)
    (f4_khipu_append_preserves_invariant es inv k newEdges hk)
\end{lstlisting}

The ASCII words \texttt{forall}, \texttt{and}, \texttt{not}, and \texttt{in}
above stand for Lean's corresponding symbols in the source. This presentation
choice keeps the arXiv listing portable; it does not claim that the normalized
excerpt itself is a compilable replacement for the repository file.

\section{Selected Lean artifact: FIFO and monotonic emit}
\label{app:f7f22}

\begin{lstlisting}[language=]
abbrev ChaskiQueue := List Nat

def chaskiEnqueue (q : ChaskiQueue) (m : Nat) : ChaskiQueue :=
  q ++ [m]

def chaskiEnqueueAll (q : ChaskiQueue) (xs : List Nat) :
    ChaskiQueue := q ++ xs

def chaskiDrain : ChaskiQueue -> List Nat
  | []      => []
  | m :: ms => m :: chaskiDrain ms

theorem f7_chaski_drain_eq (q : ChaskiQueue) :
    chaskiDrain q = q := by
  induction q with
  | nil => rfl
  | cons m ms ih => simp [chaskiDrain, ih]

theorem f7_chaski_fifo_order (msgs : List Nat) :
    chaskiDrain (chaskiEnqueueAll [] msgs) = msgs := by
  rw [f7_chaski_enqueueAll_nil, f7_chaski_drain_eq]

theorem f7_chaski_fifo_positional (msgs : List Nat) (i : Nat) :
    (chaskiDrain (chaskiEnqueueAll [] msgs))[i]? = msgs[i]? := by
  rw [f7_chaski_fifo_order]

def f22_seqLog (n : Nat) : List Nat := List.range n

theorem f22_emit_appends_length (n : Nat) :
    f22_seqLog (n + 1) = f22_seqLog n ++ [n] := by
  simp [f22_seqLog, List.range_succ]

theorem f22_khipu_emit_monotone
    (n i j : Nat) (hij : i < j) (hj : j < n) :
    (f22_seqLog n)[i] < (f22_seqLog n)[j] := by
  simp only [f22_seqLog]
  rw [List.getElem_range, List.getElem_range]
  exact hij
\end{lstlisting}

The source version of the final theorem carries the explicit bounds required by
Lean's \texttt{getElem} syntax; they are elided in the typeset excerpt but are
present at lines 254--256 of the audited file. The omission shortens the
listing, not the theorem statement analyzed in \S\ref{sec:f22}.

\section{Reviewer checklist}
\label{app:checklist}

The following checklist turns the assurance case into a repeatable review. A
release should not use the paper's certified language unless every applicable
item is answered with a commit-addressed artifact.

\begin{enumerate}
\item \textbf{Identity.} Record the 40-hex source commit, Lean toolchain, Mathlib
lock, build workflow revision, and deployed image digest.
\item \textbf{Statement parity.} Compare every theorem statement quoted by the
release with the declaration at that commit. Reject aliases that resolve to an
older vacuous theorem.
\item \textbf{Kernel check.} Run the focused module and repository build; retain
the complete log as an immutable artifact.
\item \textbf{Axiom check.} Capture per-declaration \texttt{\#print axioms};
reject \texttt{sorryAx}, undeclared project axioms, or a dependency set that
differs from the release manifest.
\item \textbf{Count parity.} Evaluate every locked-count theorem and compare the
served formula manifest with the eight expected identifiers
$\{F1,F4,F7,F11,F12,F18,F19,F22\}$.
\item \textbf{Admission mutation.} Attempt a self-reference, a forward
reference, and a mixed valid/invalid edge batch. The transaction must fail
closed without consuming a sequence number or leaving partial rows.
\item \textbf{Queue mutation.} Test two and three distinct messages, retry the
middle send, crash before and after acknowledgement, and verify the documented
ordering domain.
\item \textbf{Sequence mutation.} Race two allocators, inject a duplicate,
delete a middle row, and query through replicas/caches. The public health state
must not report a structurally valid log when any test fails.
\item \textbf{Refinement version.} Confirm that the deployed abstraction-map
version is covered by the reviewed integration tests and build attestation.
\item \textbf{Claim boundary.} Search product copy for unsupported phrases such
as ``tamper-proof,'' ``exactly once,'' ``Byzantine safe,'' or ``fully verified.''
Replace each with the narrower certified statement or attach the missing proof.
\end{enumerate}

\section{Threat-to-control traceability matrix}
\label{app:trace}

\begin{table}[h]
\centering
\small
\begin{tabular}{@{}p{0.24\linewidth}p{0.26\linewidth}p{0.38\linewidth}@{}}
\toprule
Attempt & Detecting/preventing control & Remaining obligation \\
\midrule
Circular justification & O1 plus F4 reachability proof & Adapter enforces
natural-number order atomically. \\
Self-reference on append & O2; target must be $<k$ & Batch rejection leaves no
partial receipt. \\
Message overtaking & O3 plus F7 positional equality & Broker/consumer-group
semantics refine one list order. \\
Duplicate or reused emit ID & O4 plus F22 & Concurrent allocator and persistent
store preserve the range construction. \\
Historical overwrite & Prefix relation in F22 construction & Storage snapshots
and recovery refine prefix preservation. \\
Payload substitution & Not covered structurally & Signature/hash assumption,
key identity, and verified binding path. \\
Replica disagreement & Not covered structurally & Consensus theorem or explicit
single-writer trust assumption. \\
Proof/deployment mismatch & Commit and build attestation & Deployment policy
verifies the digest and abstraction-map version. \\
\bottomrule
\end{tabular}
\caption{Traceability matrix. Blanketing all rows with ``formally verified''
would erase the distinction between proved structure and residual controls.}
\end{table}

\FloatBarrier
\footnotesize
\begin{multicols}{2}
\begin{thebibliography}{99}

\bibitem{brundage2020} M.\ Brundage et al.\ \emph{Toward Trustworthy AI
Development: Mechanisms for Supporting Verifiable Claims.} arXiv:2004.07213, 2020.
\url{https://arxiv.org/abs/2004.07213}

\bibitem{isaca2025} ISACA. \emph{The Growing Challenge of Auditing Agentic AI.}
2025. \href{https://www.isaca.org/resources/news-and-trends/industry-news/2025/the-growing-challenge-of-auditing-agentic-ai}{ISACA article}.

\bibitem{armo2025} ARMO. \emph{AI Agent Governance: the Enforceability Ladder.}
2025. \url{https://www.armosec.io/blog/ai-agent-governance/}

\bibitem{sel4} G.\ Klein et al.\ \emph{seL4: Formal Verification of an OS
Kernel.} SOSP 2009. \url{https://dl.acm.org/doi/10.1145/1629575.1629596}

\bibitem{slsa} SLSA. \emph{Supply-chain Levels for Software Artifacts, v1.0 ---
Security Levels.} \url{https://slsa.dev/spec/v1.0/levels}

\bibitem{sigstore} Sigstore. \emph{Verifying in-toto attestations with cosign.}
\url{https://docs.sigstore.dev/cosign/verifying/attestation/}

\bibitem{pbft} M.\ Castro and B.\ Liskov. \emph{Practical Byzantine Fault
Tolerance.} OSDI 1999.
\url{https://css.csail.mit.edu/6.824/2014/papers/castro-practicalbft.pdf}

\bibitem{lamport1982} L.\ Lamport, R.\ Shostak, M.\ Pease. \emph{The Byzantine
Generals Problem.} ACM TOPLAS 4(3):382--401, 1982.
\url{https://doi.org/10.1145/357172.357176}

\bibitem{velisarios} V.\ Rahli et al.\ \emph{Velisarios: Byzantine
Fault-Tolerant Protocols Powered by Coq.} ESOP 2018.
\url{https://vrahli.github.io/articles/velisarios.pdf}

\bibitem{you2019pgnn} J.\ You, R.\ Ying, J.\ Leskovec. \emph{Position-aware
Graph Neural Networks.} ICML 2019. arXiv:1906.04817.
\url{https://arxiv.org/abs/1906.04817}

\bibitem{you2020graphgym} J.\ You, R.\ Ying, J.\ Leskovec. \emph{Design Space
for Graph Neural Networks.} NeurIPS 2020. arXiv:2011.08843.
\url{https://arxiv.org/abs/2011.08843}

\bibitem{you2020graph2nn} J.\ You, J.\ Leskovec, K.\ He, S.\ Xie. \emph{Graph
Structure of Neural Networks.} ICML 2020. arXiv:2007.06559.
\url{https://arxiv.org/abs/2007.06559}

\bibitem{xu2019gin} K.\ Xu, W.\ Hu, J.\ Leskovec, S.\ Jegelka. \emph{How
Powerful are Graph Neural Networks?} ICLR 2019. arXiv:1810.00826.
\url{https://arxiv.org/abs/1810.00826}

\bibitem{mccabe2023multiphysics} M.\ McCabe et al.\ (Polymathic AI).
\emph{Multiple Physics Pretraining for Physical Surrogate Models.} 2023.
arXiv:2310.02994. \url{https://arxiv.org/abs/2310.02994}

\bibitem{mcleish2024abacus} S.\ McLeish et al.\ \emph{Transformers Can Do
Arithmetic with the Right Embeddings (Abacus Embeddings).} NeurIPS 2024.
arXiv:2405.17399. \url{https://arxiv.org/abs/2405.17399}

\bibitem{ouroboros25} S.\ P.\ Lutar Jr. \emph{Governed Post-Determinism: A
Unified Theory of Verifiable Autonomy (Ouroboros Thesis).} SZL Holdings, 2026.
Concept DOI \url{https://doi.org/10.5281/zenodo.19944926}.

\bibitem{mathlib} The mathlib Community. \emph{The Lean Mathematical Library.}
CPP 2020; L.\ de Moura, S.\ Ullrich, \emph{The Lean 4 Theorem Prover and
Programming Language}, CADE 2021.

\end{thebibliography}
\end{multicols}

\end{document}
